\chapter{Methodology}
\label{ch:methodology}

\section{Research Approach}

This thesis follows a \textbf{quantitative experimental} research strategy. The central
phenomenon—automated label template validation—is investigated by implementing a
prototype system, running controlled experiments on a labelled dataset, and measuring
performance using standard classification metrics. The study is applied in nature:
insights are derived from a real industrial problem provided by IKEA of Sweden,
and results are evaluated against operational requirements (accuracy, speed, explainability).

\section{System Overview}
\label{sec:sys_overview}

The system accepts two inputs per label validation request: (1) a structured
\acl{JSON} (JSON) metadata payload describing label content and (2) a NiceLabel
\texttt{.nlbl} label file—the engineered output produced by IKEA's label design workflow.
It produces two outputs: a binary \textbf{VALID / INVALID} decision and a structured
\textbf{validation report} detailing which elements caused any failure and where on the
label the violation occurs.

\begin{figure}[h]
\centering
\begin{tikzpicture}[node distance=0.7cm and 1.0cm]
  % Inputs
  \node[block, fill=orange!15, text width=2.6cm] (json) {JSON\\Metadata};
  \node[block, fill=orange!15, text width=2.6cm, below=of json] (nlbl) {\texttt{.nlbl}\\Label File};

  % Stage 1
  \node[block, right=1.2cm of json, yshift=-0.9cm,
        fill=blue!12, text width=3.0cm, minimum height=2.4cm]
        (stage1) {Stage 1\\[4pt]\textbf{Rendering}\\{\footnotesize NiceLabel SDK\\PNG export}};

  % Stage 2
  \node[block, right=1.0cm of stage1,
        fill=purple!10, text width=3.0cm, minimum height=2.4cm]
        (stage2) {Stage 2\\[4pt]\textbf{Pre-render}\\[2pt]\textbf{Validation}\\{\footnotesize Rule-based\\JSON checks}};

  % Stage 3
  \node[block, right=1.0cm of stage2,
        fill=teal!10, text width=3.0cm, minimum height=2.4cm]
        (stage3) {Stage 3\\[4pt]\textbf{Visual}\\[2pt]\textbf{Validation}\\{\footnotesize SSIM, ORB\\CNN, Siamese}};

  % Baseline
  \node[block, above=0.8cm of stage3, fill=green!12, text width=2.8cm]
        (baseline) {Baseline\\Template\\(PNG catalog)};

  % Output
  \node[block, right=1.0cm of stage3,
        fill=red!10, text width=2.6cm, minimum height=2.4cm]
        (out) {\textbf{Report}\\[4pt]VALID /\\INVALID\\+ BBox overlay};

  % Arrows
  \draw[arrow] (json.east) -- ++(0.5,0) |- (stage1.west);
  \draw[arrow] (nlbl.east) -- ++(0.5,0) |- (stage1.west);
  \draw[arrow] (stage1) -- (stage2);
  \draw[arrow] (stage2) -- (stage3);
  \draw[arrow] (stage3) -- (out);
  \draw[arrow] (baseline.south) -- (stage3.north);
\end{tikzpicture}
\caption{High-level system architecture showing the three validation stages.}
\label{fig:sys_overview}
\end{figure}

\section{System Architecture}
\label{sec:architecture}

The architecture follows a modular pipeline design. Each module is independently
replaceable, supporting future extension (e.g.\ swapping ORB for a detector fine-tuned
on label-specific features). The five main modules are described below and illustrated
in \cref{fig:detailed_arch}.

\begin{figure}[h]
\centering
\begin{tikzpicture}[node distance=0.55cm and 1.8cm,
  mod/.style={rectangle, rounded corners=3pt, draw=black!50,
              fill=blue!8, minimum width=3.2cm, minimum height=0.9cm,
              align=center, font=\small},
  submod/.style={rectangle, rounded corners=2pt, draw=gray!60,
                 fill=gray!8, minimum width=2.8cm, minimum height=0.7cm,
                 align=center, font=\footnotesize}
]
% Column 1: Input
\node[mod, fill=orange!15] (input) {Input Handler};
\node[submod, below=of input] (ftpmon) {FTP/API monitor};
\node[submod, below=of ftpmon] (parser) {JSON parser};

% Column 2: Template select
\node[mod, fill=purple!12, right=of input] (tmpl) {Template Selector};
\node[submod, below=of tmpl] (catalog) {Template catalog};
\node[submod, below=of catalog] (idmap) {ID $\to$ baseline PNG};

% Column 3: Renderer
\node[mod, fill=blue!12, right=of tmpl] (render) {NiceLabel Renderer};
\node[submod, below=of render] (sdk) {.NET SDK};
\node[submod, below=of sdk] (png) {PNG exporter};

% Column 4: Validator
\node[mod, fill=teal!12, right=of render] (valid) {Hybrid Validator};
\node[submod, below=of valid] (prerender) {Rule checker};
\node[submod, below=of prerender] (ssim_mod) {SSIM engine};
\node[submod, below=of ssim_mod] (orb_mod) {ORB matcher};
\node[submod, below=of orb_mod] (cnn_mod) {CNN classifier};
\node[submod, below=of cnn_mod] (siam_mod) {Siamese net};

% Column 5: Output
\node[mod, fill=red!10, right=of valid] (report) {Report Generator};
\node[submod, below=of report] (decision) {VALID/INVALID};
\node[submod, below=of decision] (bbox) {BBox overlay};
\node[submod, below=of bbox] (api) {API/Dashboard};

% Horizontal arrows
\draw[arrow] (input) -- (tmpl);
\draw[arrow] (tmpl) -- (render);
\draw[arrow] (render) -- (valid);
\draw[arrow] (valid) -- (report);
\end{tikzpicture}
\caption{Detailed modular architecture of the validation system.}
\label{fig:detailed_arch}
\end{figure}

\section{Workflow Description}
\label{sec:workflow}

\Cref{fig:workflow} shows the four-step validation workflow.

\begin{figure}[h]
\centering
\begin{tikzpicture}[node distance=1.2cm,
  step/.style={rectangle, rounded corners=5pt, draw=blue!60,
               fill=blue!10, text width=5cm, minimum height=1.2cm,
               align=center, font=\small},
  decision/.style={diamond, draw=orange!70, fill=orange!10,
                   aspect=2.5, align=center, font=\small}
]
\node[step] (s1) {\textbf{Step 1}\\Input ingestion \& template identification\\{\footnotesize Read JSON, map label ID $\to$ baseline}};
\node[step, below=of s1] (s2) {\textbf{Step 2}\\Pre-render metadata validation\\{\footnotesize Rule checks on JSON fields}};
\node[decision, below=0.7cm of s2] (d1) {Rule\\violations?};
\node[step, below=0.7cm of d1] (s3) {\textbf{Step 3}\\Rendering \& visual comparison\\{\footnotesize SDK render $\to$ PNG $\to$ SSIM/ORB/CNN/Siamese}};
\node[step, below=of s3] (s4) {\textbf{Step 4}\\Report generation\\{\footnotesize Decision + BBox overlay + API response}};

\node[step, right=2.5cm of d1, fill=red!10, draw=red!60, text width=3cm]
     (fail_early) {Early\\INVALID\\report};

\draw[arrow] (s1) -- (s2);
\draw[arrow] (s2) -- (d1);
\draw[arrow] (d1) -- node[right, font=\footnotesize]{No} (s3);
\draw[arrow] (d1) -- node[above, font=\footnotesize]{Yes} (fail_early);
\draw[arrow] (s3) -- (s4);
\end{tikzpicture}
\caption{Four-step validation workflow with early-exit on metadata failure.}
\label{fig:workflow}
\end{figure}

\subsection{Step 1 — Input Ingestion and Template Identification}

The system reads the incoming JSON payload and extracts the label identifier. A
deterministic lookup maps the identifier to the corresponding approved baseline PNG
stored in the local template catalog (approximately 20 baseline templates, with multiple
versions). Preliminary rule validation (field presence, format checks) is applied
directly to the JSON payload before any rendering occurs.

\subsection{Step 2 — Pre-Render Metadata Validation}

Rule-based checks are applied to the JSON metadata:
\begin{itemize}
  \item \textbf{Field presence}: required fields must be non-null and non-empty.
  \item \textbf{Format checks}: article numbers must match the pattern
  \texttt{\textbackslash d\{3\}.\textbackslash d\{3\}.\textbackslash d\{2\}};
  country codes must be ISO 3166-1 alpha-2; date fields must conform to YYMMDD.
  \item \textbf{Length constraints}: text fields have maximum character lengths
  defined per field per template type.
\end{itemize}

If any error-severity rule fails, an INVALID decision is returned immediately without
rendering (early exit), reducing unnecessary computation.

\subsection{Step 3 — Rendering and Visual Comparison}

The \texttt{.nlbl} file is opened by the NiceLabel .NET SDK and exported to a
300 DPI PNG. The PNG is then compared to the baseline reference using four methods
run in parallel (see \cref{sec:ml_algorithms}).

\subsection{Step 4 — Report Generation}

The output aggregates all validation findings into a structured JSON report containing:
the binary decision, per-method confidence scores, bounding-box coordinates of
localised violations, and a list of failed metadata rules with extracted values.

\section{Dataset Description}
\label{sec:dataset}

\subsection{Composition}

The dataset consists of \textbf{150 labelled label images} across 20 template types.
All valid samples are real production-approved labels; invalid samples are synthetically
generated by programmatically introducing controlled violations into valid templates.

\begin{table}[h]
\centering
\caption{Dataset composition by class and split.}
\label{tab:dataset}
\begin{tabular}{lccccc}
\toprule
\textbf{Split} & \textbf{Valid} & \textbf{Invalid} & \textbf{Total} & \textbf{\% of data} \\
\midrule
Training   & 66 & 54 & 120 & 80.0\% \\
Validation & 8  & 7  & 15  & 10.0\% \\
Test       & 8  & 7  & 15  & 10.0\% \\
\midrule
\textbf{Total} & \textbf{82} & \textbf{68} & \textbf{150} & 100\% \\
\bottomrule
\end{tabular}
\end{table}

\subsection{Synthetic Violation Types}

Seven categories of controlled violation were introduced:

\begin{table}[h]
\centering
\caption{Synthetic violation types and distribution in the invalid subset.}
\label{tab:violations}
\begin{tabular}{llc}
\toprule
\textbf{ID} & \textbf{Violation type} & \textbf{Count} \\
\midrule
V1 & Element displacement ($>$5 mm)       & 14 \\
V2 & Missing required field               & 12 \\
V3 & Font size deviation ($>$15\%)        & 10 \\
V4 & Barcode zone overlap                 & 9  \\
V5 & Incorrect article number format      & 9  \\
V6 & Colour value mismatch                & 8  \\
V7 & Extra / unexpected element           & 6  \\
\midrule
\textbf{Total}                            && 68 \\
\bottomrule
\end{tabular}
\end{table}

\subsection{Data Split Rationale}

The 80/10/10 split was chosen over a larger test set due to the small overall dataset
size. The same split proportions were maintained within each template type to avoid
type-level imbalance. Stratified splitting was applied at the violation-type level to
ensure all seven violation categories are represented in both validation and test sets.

\section{Validation Strategy}
\label{sec:validation_strategy}

The validation strategy is \textbf{hybrid}: four complementary methods are applied to
the rendered label, and their decisions are combined through a weighted majority vote.

\begin{figure}[h]
\centering
\begin{tikzpicture}[node distance=0.7cm and 1.4cm,
  mbox/.style={rectangle, rounded corners=3pt, draw=black!50, fill=gray!8,
               text width=3.0cm, minimum height=1.1cm, align=center, font=\small}
]
\node[mbox, fill=blue!12]  (ssim_n) {SSIM\\baseline};
\node[mbox, fill=purple!12, right=of ssim_n]   (orb_n)  {ORB feature\\matching};
\node[mbox, fill=teal!12,   right=of orb_n]    (cnn_n)  {ResNet-50\\classifier};
\node[mbox, fill=orange!12, right=of cnn_n]    (siam_n) {Siamese\\network};

\node[mbox, fill=red!10, draw=red!60, below=1.2cm of orb_n, xshift=2.5cm,
      text width=4cm, minimum height=1.2cm]
     (vote) {Weighted majority vote\\{\footnotesize w = [0.15, 0.25, 0.35, 0.25]}};

\node[mbox, fill=green!10, draw=green!60!black, right=1.2cm of vote,
      text width=2.5cm]
     (dec) {VALID /\\INVALID};

\draw[arrow] (ssim_n.south) -- ++(0,-0.5) -| (vote.north);
\draw[arrow] (orb_n.south)  -- ++(0,-0.3) -| (vote.north);
\draw[arrow] (cnn_n.south)  -- ++(0,-0.3) -| (vote.north);
\draw[arrow] (siam_n.south) -- ++(0,-0.5) -| (vote.north);
\draw[arrow] (vote) -- (dec);
\end{tikzpicture}
\caption{Hybrid validation ensemble: four methods combined by weighted majority vote.}
\label{fig:ensemble}
\end{figure}

The weights $w = [0.15, 0.25, 0.35, 0.25]$ for SSIM, ORB, CNN, and Siamese
respectively were determined empirically on the validation set using a grid search over
integer weight combinations that sum to 1.00.

\section{Machine Learning Algorithms}
\label{sec:ml_algorithms}

\subsection{Method 1 — SSIM Baseline}

The Structural Similarity Index Measure~\cite{wang2004ssim}:

\begin{equation}
\text{SSIM}(x, y) = \frac{(2\mu_x\mu_y + c_1)(2\sigma_{xy} + c_2)}
                         {(\mu_x^2 + \mu_y^2 + c_1)(\sigma_x^2 + \sigma_y^2 + c_2)}
\label{eq:ssim}
\end{equation}

is computed between the generated label PNG and the baseline reference PNG after
aligning both images to the same resolution. A sliding-window SSIM map with window
size $11 \times 11$ pixels is used; regions where local SSIM falls below a threshold
$\tau_\text{SSIM} = 0.82$ are flagged as violated areas. Bounding boxes are derived by
morphologically dilating the flagged pixels and extracting connected components.

\subsection{Method 2 — ORB Feature Matching}

Oriented FAST and Rotated BRIEF (ORB)~\cite{rublee2011orb} keypoints are extracted
from both the generated label and the baseline reference. Keypoint descriptors are
matched using a Hamming-distance Brute-Force matcher with ratio test threshold~0.75
(Lowe's ratio test criterion). A homography is estimated from matched keypoints using
RANSAC; the inlier ratio $r_\text{in}$ is used as the similarity score. An INVALID
decision is issued when $r_\text{in} < 0.60$. Additionally, unmatched keypoints in the
generated label that have no corresponding counterpart within 25 pixels in the reference
are treated as potential anomalous regions.

\subsection{Method 3 — ResNet-50 Binary Classifier}
\label{ssec:resnet}

A ResNet-50 backbone~\cite{he2016resnet} pre-trained on ImageNet is fine-tuned as a
binary classifier (VALID / INVALID). The classification head replaces the original
1\,000-class softmax with a two-layer MLP:

\begin{equation}
\hat{y} = \sigma\!\left(W_2\,\text{ReLU}(W_1\,z + b_1) + b_2\right)
\label{eq:cls_head}
\end{equation}

where $z \in \mathbb{R}^{2048}$ is the global average-pooled ResNet-50 feature vector,
$W_1 \in \mathbb{R}^{512 \times 2048}$, $W_2 \in \mathbb{R}^{1 \times 512}$, and
$\sigma$ is the sigmoid function. The network is trained end-to-end with binary
cross-entropy loss:

\begin{equation}
\mathcal{L}_\text{BCE} = -\frac{1}{N}\sum_{i=1}^{N}
  \left[y_i \log \hat{y}_i + (1-y_i)\log(1-\hat{y}_i)\right]
\label{eq:bce}
\end{equation}

\paragraph{Training details.}

\begin{table}[h]
\centering
\caption{ResNet-50 classifier training hyperparameters.}
\label{tab:resnet_hparams}
\begin{tabular}{ll}
\toprule
\textbf{Hyperparameter} & \textbf{Value} \\
\midrule
Backbone         & ResNet-50 (ImageNet pre-trained) \\
Fine-tuning      & Full network (unfrozen after epoch~5) \\
Optimiser        & Adam ($\beta_1 = 0.9$, $\beta_2 = 0.999$) \\
Learning rate    & $1\times10^{-4}$ (backbone), $1\times10^{-3}$ (head) \\
LR schedule      & CosineAnnealingLR, $T_\text{max} = 50$ \\
Batch size       & 16 \\
Epochs           & 50 (early stopping patience = 8) \\
Input resolution & $224 \times 224$ px (resized from 300 DPI PNG) \\
Augmentation     & Random horizontal flip, $\pm$10° rotation,\\
                 & colour jitter (brightness/contrast $\pm$0.2) \\
Class weights    & Inverse frequency ($w_\text{valid}=0.83$, $w_\text{invalid}=1.00$) \\
\bottomrule
\end{tabular}
\end{table}

\Cref{fig:resnet_arch} illustrates the modified ResNet-50 architecture.

\begin{figure}[h]
\centering
\begin{tikzpicture}[node distance=0.5cm and 0.8cm,
  layer/.style={rectangle, rounded corners=2pt, draw=black!50,
                fill=blue!10, minimum width=2.2cm, minimum height=0.8cm,
                align=center, font=\scriptsize}
]
\node[layer, fill=orange!15] (inp) {Input\\$224\times224\times3$};
\node[layer, right=of inp] (conv1) {Conv1 + BN\\+ ReLU + Pool};
\node[layer, right=of conv1] (res1) {ResBlock\\$\times3$ (conv2)};
\node[layer, right=of res1]  (res2) {ResBlock\\$\times4$ (conv3)};
\node[layer, right=of res2]  (res3) {ResBlock\\$\times6$ (conv4)};
\node[layer, right=of res3]  (res4) {ResBlock\\$\times3$ (conv5)};
\node[layer, below=0.7cm of res4, fill=gray!15] (gap)  {Global Avg Pool\\$\mathbb{R}^{2048}$};
\node[layer, left=of gap, fill=teal!15]  (fc1)  {FC 2048$\to$512\\ReLU};
\node[layer, left=of fc1, fill=teal!15]  (fc2)  {FC 512$\to$1\\Sigmoid};
\node[layer, left=of fc2, fill=red!15]   (out)  {VALID /\\INVALID};

\draw[arrow] (inp)  -- (conv1);
\draw[arrow] (conv1)-- (res1);
\draw[arrow] (res1) -- (res2);
\draw[arrow] (res2) -- (res3);
\draw[arrow] (res3) -- (res4);
\draw[arrow] (res4) -- (gap);
\draw[arrow] (gap)  -- (fc1);
\draw[arrow] (fc1)  -- (fc2);
\draw[arrow] (fc2)  -- (out);
\end{tikzpicture}
\caption{Modified ResNet-50 architecture with binary classification head.}
\label{fig:resnet_arch}
\end{figure}

\subsection{Method 4 — Siamese Network}
\label{ssec:siamese}

A Siamese network~\cite{sung2018relation} learns a pairwise similarity function
$f: (\mathbf{x}_\text{gen}, \mathbf{x}_\text{ref}) \to [0, 1]$ directly from
(generated label, reference label) pairs. Two ResNet-18 sub-networks with shared
weights encode each image independently into a $d = 512$ embedding:

\begin{equation}
z_i = \text{ResNet18}(\mathbf{x}_i),\quad z_i \in \mathbb{R}^{512}
\label{eq:encoder}
\end{equation}

A relation module computes:
\begin{equation}
s = \sigma\!\left(r\!\left([z_\text{gen} \,\Vert\, z_\text{ref}]\right)\right)
\label{eq:relation}
\end{equation}

where $[\,\cdot\,\Vert\,\cdot\,]$ is concatenation, $r(\cdot)$ is a two-layer MLP
($\mathbb{R}^{1024} \to \mathbb{R}^{256} \to \mathbb{R}^{1}$), and $s \in [0,1]$ is
the similarity score ($s < 0.5 \Rightarrow$ INVALID). The model is trained with
contrastive loss:

\begin{equation}
\mathcal{L}_\text{cont} = \frac{1}{2N}\sum_{i=1}^N
  \left[y_i \cdot d_i^2 + (1-y_i) \cdot \max(m - d_i, 0)^2\right]
\label{eq:contrastive}
\end{equation}

where $d_i = \|z_{\text{gen},i} - z_{\text{ref},i}\|_2$, $y_i \in \{0,1\}$
is the similarity label (1 = same class), and $m = 1.0$ is the margin.

\begin{table}[h]
\centering
\caption{Siamese network training hyperparameters.}
\label{tab:siamese_hparams}
\begin{tabular}{ll}
\toprule
\textbf{Hyperparameter} & \textbf{Value} \\
\midrule
Backbone (shared)   & ResNet-18 (ImageNet pre-trained, unfrozen from epoch 3) \\
Relation MLP        & 1024 $\to$ 256 $\to$ 1 (ReLU, Sigmoid) \\
Optimiser           & Adam ($\beta_1 = 0.9$, $\beta_2 = 0.999$) \\
Learning rate       & $5\times10^{-5}$ \\
Margin $m$          & 1.0 \\
Pair sampling       & 3 positive + 3 negative pairs per anchor \\
Batch size          & 12 pairs \\
Epochs              & 40 (early stopping patience = 6) \\
Input resolution    & $224 \times 224$ px \\
\bottomrule
\end{tabular}
\end{table}

\Cref{fig:siamese_arch} illustrates the Siamese architecture.

\begin{figure}[h]
\centering
\begin{tikzpicture}[node distance=0.6cm and 1.2cm,
  enc/.style={rectangle, rounded corners=3pt, draw=blue!60, fill=blue!12,
              text width=2.4cm, minimum height=0.9cm, align=center, font=\small},
  rel/.style={rectangle, rounded corners=3pt, draw=teal!60, fill=teal!10,
              text width=2.6cm, minimum height=0.9cm, align=center, font=\small}
]
% Generated branch
\node[enc] (gen_in) {Generated\\label PNG};
\node[enc, below=0.5cm of gen_in, fill=blue!20] (gen_enc) {ResNet-18\\encoder};
\node[rel, right=1.2cm of gen_enc, yshift=0.7cm] (concat) {Concat\\$z_\text{gen}\,||\,z_\text{ref}$};

% Reference branch
\node[enc, below=1.5cm of gen_in] (ref_in) {Reference\\baseline PNG};
\node[enc, below=0.5cm of ref_in, fill=blue!20] (ref_enc) {ResNet-18\\encoder\\(shared weights)};

% Relation
\node[rel, right=1.2cm of concat] (rel_mod) {Relation\\MLP};
\node[enc, right=1.2cm of rel_mod, fill=red!10, draw=red!60] (score) {Similarity\\score $s \in [0,1]$};

\draw[arrow] (gen_in) -- (gen_enc);
\draw[arrow] (ref_in) -- (ref_enc);
\draw[arrow] (gen_enc.east) -- ++(0.3,0) |- (concat.west);
\draw[arrow] (ref_enc.east) -- ++(0.3,0) |- (concat.west);
\draw[arrow] (concat) -- (rel_mod);
\draw[arrow] (rel_mod) -- (score);
\end{tikzpicture}
\caption{Siamese network architecture for pairwise template similarity scoring.}
\label{fig:siamese_arch}
\end{figure}

\section{Evaluation Method}
\label{sec:evaluation_method}

\subsection{Metrics}

Each method is evaluated on the test set (15 labels) using:

\begin{itemize}
  \item \textbf{Accuracy}: $\text{Acc} = (\text{TP}+\text{TN})/N$
  \item \textbf{Precision}: $P = \text{TP}/(\text{TP}+\text{FP})$
  \item \textbf{Recall}: $R = \text{TP}/(\text{TP}+\text{FN})$
  \item \textbf{F1-score}: $F1 = 2PR/(P+R)$
  \item \textbf{Violation localisation}: IoU between predicted and annotated bounding
  boxes; a prediction is correct if IoU~$>$~0.5.
  \item \textbf{Processing time}: wall-clock time from PNG ingestion to report
  generation, averaged over the test set.
\end{itemize}

\subsection{Baseline}

A pixel-difference baseline computes the $L_1$ mean absolute pixel difference between
the generated and reference PNG. Labels with mean difference exceeding 12 intensity
units (on a 0–255 scale) are classified as INVALID. This baseline is included to
confirm that classical methods are insufficient on their own.

\section{Method Justification}
\label{sec:justification}

The choice of a hybrid approach is motivated by the literature's consistent finding that
no single comparison method provides sufficient coverage across all violation
types~\cite{zhang2018lpips, roth2022patchcore}. SSIM is fast and interpretable but
sensitive to minor rendering variations; ORB is robust to slight positional shifts but
struggles on uniform layouts with few distinctive keypoints; the ResNet-50 classifier
captures global structural patterns missed by local metrics; the Siamese network enables
few-shot comparison directly from reference pairs without retraining.

Reference-based validation was chosen because IKEA of Sweden's workflow provides
an authoritative approved baseline for every label type. This makes the problem
well-suited to deterministic comparison~\cite{bergmann2019mvtec}. The deliberate
exclusion of synthetic negative data for training the rule-based component reflects
the practical reality that failed labels are not archived in the production system.
